\documentclass[11pt]{article}

\usepackage[T1]{fontenc}
\usepackage{lmodern}
\usepackage{amsmath,amssymb,amsthm}
\usepackage{booktabs}
\usepackage[margin=1in]{geometry}
\usepackage{microtype}
\usepackage[hidelinks]{hyperref}
\hypersetup{
  pdftitle={A Counterexample to Conjecture 1 on Pseudo-Random Array Codes},
  pdfauthor={Cathal Ryan Hynes},
  pdfsubject={Pseudo-random array codes and de Bruijn array codes},
  pdfkeywords={pseudo-random array code, PRAC, de Bruijn array code,
    finite fields, linear recurrence}
}

\newtheorem{theorem}{Theorem}
\newtheorem{lemma}[theorem]{Lemma}
\newtheorem{proposition}[theorem]{Proposition}
\newtheorem{corollary}[theorem]{Corollary}

\newcommand{\F}{\mathbf F}
\newcommand{\Tr}{\operatorname{Tr}}
\newcommand{\ord}{\operatorname{ord}}

\title{A Counterexample to Conjecture 1 on\\
Pseudo-Random Array Codes}
\author{Cathal Ryan Hynes\\
\small P.I.G.M.I.E. Ltd. (PolieBotics), Ireland}
\date{21 July 2026}

\begin{document}
\maketitle

\begin{abstract}
\noindent
Conjecture 1 of Blackburn, Chee, Etzion and Lao,
\emph{On de Bruijn Array Codes, Part II: Pseudo-Random Array Codes},
asserts a product-closure property for certain pseudo-random array codes
under the side condition \(n_1<r_1<2n_1\).
We give a counterexample with
\[
 (r_1,r_2;n_1,n_2)=(5,117;3,4).
\]
Two explicit irreducible binary polynomials of degree \(12\) and exponent
\(585\) each yield a \((5,117;3,4)\)-PRAC, while their product does not
yield a \((5,117;3,8)\)-PRAC.
The failure has a short finite-field certificate: a nonzero binary
polynomial supported on the \(3\) by \(8\) rectangle vanishes at both
evaluation pairs.
A separately written recurrence-and-folding implementation exhaustively
verifies both hypotheses and exhibits a concrete repeated window.
We also prove the conjectured product closure for the subfamily
\(\ord_{r_1}(2)=n_1\), for every number of factors.
\end{abstract}

\section{Introduction}

Etzion developed de Bruijn array codes and their shortened variants in
Part I~\cite{etzion-part1}.  Blackburn, Chee, Etzion and Lao developed the
linear theory in Part II~\cite{blackburn-part2}.  We use the definitions and
folding convention of the version of record of Part II; where its
numbering differs from the arXiv preprint, print numbering is used,
with the arXiv numbers catalogued in \nolinkurl{NUMBERING_MAP.md} of
the accompanying repository.

An \((r_1,r_2;n_1,n_2)\) pseudo-random array code, or PRAC, is a collection
of binary doubly periodic \(r_1\) by \(r_2\) arrays in which every nonzero
binary \(n_1\) by \(n_2\) matrix occurs exactly once as a toroidal window,
and which has the shift-and-add property.  When
\(\gcd(r_1,r_2)=1\), a period-\(r_1r_2\) sequence is folded by placing
symbol \(s_t\) at the unique coordinate \((i,j)\) satisfying
\[
 t\equiv i\pmod {r_1},\qquad t\equiv j\pmod {r_2}.
\]
Throughout, claims are proved by the argument supplied in the text or
verified by the recorded computations that accompany this manuscript,
as indicated in context.

The current form of Conjecture 1 in Part II states the following.
Suppose
\[
 n_1<r_1<2n_1
\]
and \(f_1,\ldots,f_\ell\) are irreducible binary polynomials of degree
\(n_1n_2\) and exponent \(r_1r_2\), each of whose folded sequences forms an
\((r_1,r_2;n_1,n_2)\)-PRAC.  Then the conjecture asserts that the product of
any \(k\) of the factors yields an
\((r_1,r_2;n_1,kn_2)\)-PRAC.

The counterexample below satisfies the strict side condition.  It therefore
resolves Conjecture 1 as stated, but makes no claim about a complete
classification of product closure.

\section{The evaluation criterion}

We first put Theorem 6 of Part II (Theorem 47 in the arXiv version)
into a form suited to a short certificate.  Let
\[
 h=n_1,\qquad w=n_2,\qquad n=hw,\qquad e=r_1r_2.
\]
Choose integers \(\mu,\nu\) with
\(\mu r_1+\nu r_2=1\).  If \(\alpha_u\) is a root of \(f_u\), define
\[
 \beta_u=\alpha_u^{\nu r_2},\qquad
 \gamma_u=\alpha_u^{\mu r_1}.
\]
Thus \(\alpha_u=\beta_u\gamma_u\), while \(\beta_u\) and \(\gamma_u\)
have orders \(r_1\) and \(r_2\), respectively.

\begin{lemma}[Evaluation form of Theorem 6]
\label{lem:evaluation}
For \(k\) selected factors, their product folding gives an
\((r_1,r_2;h,kw)\)-PRAC if and only if the simultaneous evaluation map
\[
\begin{split}
 \operatorname{ev}\colon
 \{P\in\F_2[X,Y]:\deg_XP<h,\ \deg_YP<kw\}
 &\longrightarrow \bigoplus_{u=1}^k\F_{2^n},\\
 P&\longmapsto
 \bigl(P(\beta_u,\gamma_u)\bigr)_{u=1}^k
\end{split}
\]
is injective.
\end{lemma}

\begin{proof}
For each pair \(0\leq i<h\), \(0\leq j<kw\), consider
\[
 v_{ij}=(\beta_1^i\gamma_1^j,\ldots,
          \beta_k^i\gamma_k^j)
 \in\bigoplus_{u=1}^k\F_{2^n}.
\]
Taking \(t_u=\Tr\), as Theorem 6 permits, the \(u\)-th block of
row \((i,j)\) in the matrix of Theorem 6 consists of
the absolute trace pairings of \(\beta_u^i\gamma_u^j\) with the basis
\(1,\alpha_u,\ldots,\alpha_u^{n-1}\).  The trace pairing is nondegenerate.
Consequently the matrix in Theorem 6 has full rank exactly when the
\(kn\) vectors \(v_{ij}\) are linearly independent over \(\F_2\).  This is
equivalent to injectivity of the displayed evaluation map.
\end{proof}

For one factor and a width-\(w\) window, Lemma~\ref{lem:evaluation} is the
linear-independence criterion in Corollary 15 of Part II
(arXiv Corollary 48).  We use the
index range \(0\leq j<w\), as in the proof of that corollary.

\section{The counterexample}

\begin{theorem}
\label{thm:counterexample}
Let
\[
\begin{aligned}
 f_1(X)&=X^{12}+X^{10}+X^6+X^5+X^3+X+1,\\
 f_2(X)&=X^{12}+X^9+X^8+X^4+X^3+X+1.
\end{aligned}
\]
Both \(f_1\) and \(f_2\) are irreducible of degree \(12\) and exponent
\(585=5\cdot117\).  Folding the sequences generated by either factor
yields a \((5,117;3,4)\)-PRAC, but folding the sequences generated by
\(f_1f_2\) does not yield a \((5,117;3,8)\)-PRAC.
\end{theorem}

Since \(3<5<6\), Theorem~\ref{thm:counterexample} is a counterexample to
Conjecture 1.

\subsection{Irreducibility and exponent}

We encode a binary polynomial by the integer whose bit \(i\) is the
coefficient of \(X^i\).  Thus \(f_1=\mathtt{0x146b}\) and
\(f_2=\mathtt{0x131b}\).  Rabin's irreducibility test gives
\[
 X^{2^{12}}\equiv X\pmod {f_u},
\]
and
\[
\gcd(f_u,X^{2^6}+X)=\gcd(f_u,X^{2^4}+X)=1
\quad (u=1,2).
\]
The exact remainders needed to check the order are
\[
\begin{array}{c|cc}
 &\bmod f_1&\bmod f_2\\ \hline
X^{585}&\mathtt{0x001}&\mathtt{0x001}\\
X^{195}&\mathtt{0xab8}&\mathtt{0x520}\\
X^{117}&\mathtt{0x4a8}&\mathtt{0x7ce}\\
X^{45} &\mathtt{0x19c}&\mathtt{0x73d}
\end{array}
\]
The first row shows \(X^{585}\equiv1\pmod{f_u}\), while the remaining
rows show \(X^{585/p}\not\equiv1\pmod{f_u}\) for each prime divisor
\(p\in\{3,5,13\}\) of \(585\).  Together these prove that each
residue class of \(X\) has exact order \(585\).

\subsection{The two single-factor tests}

Work in
\[
 K=\F_2[a]/(f_1(a)).
\]
Direct polynomial reduction shows that \(a_2=a^{46}\) is a root of \(f_2\).
We take \(\alpha_1=a\) and \(\alpha_2=a_2\).
Use
\[
 47\cdot5-2\cdot117=1
\]
and define
\[
 b=a^{351}=a_2^{351},\qquad
 c_1=a^{235},\qquad c_2=a_2^{235}=a^{280}.
\]
The element \(b\) has order \(5\), so
\[
 F=\F_2(b)\cong
 \F_2[B]/(B^4+B^3+B^2+B+1)=\F_{16}.
\]
Both \(c_1\) and \(c_2\) have degree three over \(F\).  Their distinct
minimal polynomials over \(F\) are
\[
\begin{aligned}
 m_1(Y)&=Y^3+(1+b+b^2+b^3)Y^2+b^2Y+(b^2+b^3),\\
 m_2(Y)&=Y^3+(1+b+b^2+b^3)Y^2
              +(1+b+b^3)Y+(b^2+b^3).
\end{aligned}
\]

Put
\[
 U=\operatorname{span}_{\F_2}\{1,b,b^2\}\subset F.
\]
The binary \(3\) by \(4\) rectangle evaluates at \(X=b\) to the
\(\F_2\)-space \(U[Y]_{<4}\).
If a nonzero member of this space vanished at \(c_u\), it would have the
form \(q\,m_u\), where
\[
 q=q_0+q_1b+q_2b^2\in U\setminus\{0\}.
\]
Requiring the \(b^3\)-coordinate of all four coefficients to vanish gives,
for either \(m_u\),
\[
 q_0=0,\qquad q_0+q_2=0.
\]
The coefficient of \(Y\) then gives \(q_1+q_2=0\) for \(m_1\), or
\(q_0+q_1+q_2=0\) for \(m_2\).  Hence \(q_0=q_1=q_2=0\), a contradiction.
Lemma~\ref{lem:evaluation} therefore proves both single-factor PRAC claims.

\subsection{A rectangle relation for the product}

Multiplication in \(F[Y]\) gives
\[
\begin{aligned}
 M(Y):=m_1(Y)m_2(Y)
 &=Y^6+(1+b+b^2)Y^4+b^3Y^3\\
 &\quad +(1+b^2+b^3)Y^2+(b+b^2)Y+(1+b^2+b^3).
\end{aligned}
\]
Although \(M\) has coefficients outside \(U\), multiplication by \(Y+b\)
cancels all \(b^3\)-coordinates:
\[
\begin{aligned}
 A(Y):=(Y+b)M(Y)
 &=Y^7+bY^6+(1+b+b^2)Y^5+(b+b^2)Y^4\\
 &\quad+bY^3+(1+b)Y^2+Y+(1+b^2).
\end{aligned}
\]
Thus \(A\in U[Y]_{<8}\setminus\{0\}\), and
\[
 A(c_1)=A(c_2)=0.
\]
Expressing each coefficient in the basis \(1,b,b^2\) produces the binary
rectangle polynomial
\[
\boxed{
\begin{aligned}
P(X,Y)={}&1+Y+Y^2+Y^5+Y^7\\
 &+X(Y^2+Y^3+Y^4+Y^5+Y^6)\\
 &+X^2(1+Y^4+Y^5).
\end{aligned}}
\]
It is nonzero, \(\deg_XP<3\), \(\deg_YP<8\), and
\[
 P(b,c_1)=P(b,c_2)=0.
\]
Lemma~\ref{lem:evaluation} now proves that the product fails.

There is also a basis-free univariate certificate.  For a root
\(\alpha_u\) of either factor,
\[
 \beta_u=\alpha_u^{351},\qquad
 \gamma_u=\alpha_u^{235}.
\]
Reducing exponents modulo \(585\), the rectangle relation becomes
\[
\begin{aligned}
 Q(Z)&=P(Z^{351},Z^{235})\pmod {Z^{585}-1}\\
 &=1+Z^5+Z^6+Z^{117}+Z^{121}+Z^{122}
   +Z^{235}+Z^{236}\\
 &\quad+Z^{356}+Z^{470}+Z^{471}+Z^{472}+Z^{475}.
\end{aligned}
\]
Direct binary polynomial division gives
\[
 Q(Z)\equiv0\pmod {f_1(Z)},\qquad
 Q(Z)\equiv0\pmod {f_2(Z)}.
\]
This sparse identity is an independently checkable certificate of the
rank failure; it uses neither a field basis nor a trace convention.

\section{Exhaustive verification by a separate implementation}

A separate implementation generated sequences directly from the binary
linear recurrences, folded them using only the CRT rule, and exhaustively
enumerated all windows.  It was written with knowledge of the claimed
witness, but reuses no source code from the algebraic path and shares
no finite-field, trace, determinant, or evaluation machinery with the
proof above.

For each single factor, the checker enumerated
\[
 \frac{2^{12}-1}{585}=7
\]
shift-orbit representatives and scanned all \(7\cdot585=4095\) windows.
Every nonzero \(3\) by \(4\) binary matrix occurred exactly once, and no
zero window occurred.

For the product, the checker enumerated
\[
 \frac{2^{24}-1}{585}=28{,}679
\]
shift-orbit representatives and scanned
\[
 28{,}679\cdot585=16{,}777{,}215=2^{24}-1
\]
windows.  The results were:
\[
\begin{array}{lr}
\toprule
\text{quantity}&\text{count}\\
\midrule
\text{distinct nonzero windows observed}&8{,}388{,}607\\
\text{missing nonzero windows}&8{,}388{,}608\\
\text{duplicate nonzero occurrences}&8{,}388{,}607\\
\text{zero-window occurrences}&1\\
\bottomrule
\end{array}
\]

With rows written left to right, the window
\[
\begin{array}{c}
11110001\\
00010011\\
11100011
\end{array}
\]
occurs in two distinct shift orbits.  Using zero-based coordinates, its
two frozen occurrences are
\[
\begin{array}{c|cc}
\text{representative state}&\text{top}&\text{left}\\ \hline
1&1&40\\
8&4&28
\end{array}
\]
The first missing nonzero window is
\[
\begin{array}{c}
00100000\\
00000000\\
00000000
\end{array},
\]
and an all-zero window occurs in representative state \(10781\) at
\((\text{top},\text{left})=(4,14)\).

The frozen JSON record has SHA-256
\[
\mathtt{ba098b857e2aeeccef726e3440736dd25969732cd6c0bac0cb984a43efaa3b97},
\]
and the checker has SHA-256
\[
\mathtt{ed84442dc519ff07fc24c27f00d4b91366011c89a24e70562349351875ba23f0}.
\]
The run used Python 3.12.13 on Linux and no randomness.  The complete
record and checker are supplied as
\nolinkurl{implementations/bruteforce/counterexample_frozen.json} and
\nolinkurl{implementations/bruteforce/bruteforce.py}.

\section{Why the side condition was not enough}

The counterexample suggests a useful correction to the structural picture.
We retain \(h,w,r_1,r_2,n\) from Section 2 and write
\[
 d=\ord_{r_1}(2).
\]

\begin{lemma}[Normalisation]
\label{lem:normalisation}
If one single factor passes and \(h<r_1<2h\), then
\[
 d=\varphi(r_1),\qquad h\leq d<2h.
\]
Consequently, after replacing each selected root by a binary Frobenius
conjugate, all its \(r_1\)-components may be chosen equal:
\[
 \beta_1=\cdots=\beta_k=b.
\]
\end{lemma}

\begin{proof}
The single-factor basis contains
\(1,\beta,\ldots,\beta^{h-1}\), so these powers are independent and
\[
 d=[\F_2(\beta):\F_2]\geq h.
\]
Since \(d\mid\varphi(r_1)\) and
\(\varphi(r_1)\leq r_1-1<2h\), the positive integer
\(\varphi(r_1)/d\) is less than two and hence equals one.
Thus all primitive \(r_1\)-th roots form one Frobenius orbit, proving the
normalisation claim.
\end{proof}

Let
\[
 F=\F_2(b)=\F_{2^d},\qquad
 U=\operatorname{span}_{\F_2}\{1,b,\ldots,b^{h-1}\}.
\]
Let \(g_u(Y)\in F[Y]\) be the minimal polynomial of \(\gamma_u\) over
\(F\), and put \(m=n/d\).  Since
\(\alpha_u=b\gamma_u\) and \(K=\F_2(\alpha_u)\), one has
\(F(\gamma_u)=K\); hence \(\deg g_u=m\).
Distinct binary factors give distinct \(g_u\): equality would make
\(\gamma_v\) an \(F\)-Frobenius conjugate of \(\gamma_u\), and therefore
\(\alpha_v\) a binary Frobenius conjugate of \(\alpha_u\).

The evaluation criterion now becomes the exact coefficient-space condition
\[
 U[Y]_{<kw}\cap
 \left(\prod_{u=1}^k g_u(Y)\right)=\{0\},
\]
where the ideal is taken in \(F[Y]\).  The individual hypotheses are
\[
 U[Y]_{<w}\cap(g_u)=\{0\}.
\]
If \(d=h\), then \(U=F\), and ordinary degree counting proves product
closure.  If \(d>h\), then \(U\) is only a proper additive subspace of
\(F\), and \(m=hw/d<w\).  A low-degree multiple of
\(\prod g_u\) can therefore have all its coefficients in \(U\).
In Theorem~\ref{thm:counterexample}, \(d=4\), \(h=3\), \(m=3\), \(w=4\),
and the multiplier is exactly \(Y+b\).

\begin{theorem}[A product-closed subfamily]
\label{thm:subfamily}
Let \(h,w,r,s\) be positive integers such that
\[
 \gcd(r,s)=1,\qquad h<r<2h,\qquad
 \ord_{rs}(2)=hw,\qquad \ord_r(2)=h.
\]
For every collection of \(k\) distinct irreducible binary polynomials of
degree \(hw\) and exponent \(rs\), folding the sequences generated by their
product yields an \((r,s;h,kw)\)-PRAC.
\end{theorem}

\begin{proof}
Since \(\ord_r(2)=h\), while
\(\ord_r(2)\mid\varphi(r)\leq r-1<2h\), we have
\(\ord_r(2)=\varphi(r)\).  Thus all primitive order-\(r\) elements form one
binary Frobenius orbit, and we may choose the roots so their order-\(r\)
components all equal \(b\).  Here
\[
 F=\F_2(b)=\F_{2^h}
\]
and \(1,b,\ldots,b^{h-1}\) is a basis of \(F\), so \(U=F\).
For each selected root, its order-\(s\) component \(\gamma_u\) has degree
\[
 [\F_{2^{hw}}:F]=w
\]
over \(F\).  The corresponding minimal polynomials \(g_u\in F[Y]\) are
distinct, by the argument preceding the theorem.  A binary rectangle
relation would give a polynomial \(A\in F[Y]\) of degree less than \(kw\)
that vanishes at every \(\gamma_u\).  It would therefore be divisible by
the product of the \(k\) distinct degree-\(w\) polynomials \(g_u\), which
has degree \(kw\).  Hence \(A=0\), and then all binary rectangle
coefficients are zero.  Lemma~\ref{lem:evaluation} completes the proof.
\end{proof}

Theorem~\ref{thm:subfamily} covers the sufficient-condition family in
Theorem 5 of Part II (Theorem 30 in the arXiv version) under the
conjecture's side condition.  It includes
the motivating \((5,17;4,2)\) example and the
\((3,7;2,3)\) exponent-\(21\) example.

Lemma 17 of Part II (Lemma 44 in the arXiv version) does not prove
this result.  That lemma concerns the
paper's \(\vee\)-construction, with separate degree-\(n_1\),
exponent-\(r_1\) and degree-\(n_2\), exponent-\(r_2\) factors.
Conjecture 1 multiplies arbitrary degree-\(n_1n_2\),
exponent-\(r_1r_2\) factors.  In the counterexample there is no
degree-three subfield containing \(b\), so the required decomposition is
absent.

\section{Conclusion}

Theorem~\ref{thm:counterexample} settles Conjecture 1 in the negative.
The sparse polynomial \(P\) is a compact certificate: it is a
nonzero member of the proposed \(3\) by \(8\) window space and vanishes at
both selected evaluation pairs.  The separately written exhaustive recurrence computation confirms
both single-factor hypotheses and the product failure.

Theorem~\ref{thm:subfamily} of this note (unrelated to Theorem 4 of
the published source paper, the join-construction result) is a
sufficient condition for product closure: it covers the exponent-85
motivating family and identifies a broad sufficient regime.  Outside
that regime, the note gives a coefficient-space intersection
criterion, not a complete classification: product closure is
equivalent to
\[
 U[Y]_{<kw}\cap(\prod g_u)=\{0\}.
\]
This formulation may be useful for classifying the remaining parameters.

\subsection*{AI-assistance disclosure}

This research was produced by an AI research pipeline directed by
Cathal Ryan Hynes at P.I.G.M.I.E. Ltd., Ireland.  From a research
brief written by Anthropic's Claude (Fable), OpenAI's Sol Ultra, running
in the ChatGPT web interface, found the counterexample and the
sufficient-condition theorem; Fable separately rebuilt the entire
verification, reusing no source code, and matched every count; a separate Sol Ultra
session contributed two further audit rounds.  Cathal Ryan Hynes set the
target and the verification standard, checked the certificate himself,
and stands behind this release.  The exhaustive recurrence checker, the separate algebra
implementation, the third finite-field implementation, the audit
reimplementation, and the frozen outputs and hashes accompany this
manuscript.  No claim is made beyond resolution of Conjecture 1 as
stated in the version of record and the explicit sufficient-condition
theorem proved here.

\begin{thebibliography}{9}

\bibitem{blackburn-part2}
S.~R. Blackburn, Y.~M. Chee, T.~Etzion, and H.~Lao,
\newblock On de Bruijn array codes, Part II: Pseudo-random array codes,
\newblock \emph{IEEE Transactions on Information Theory},
72(8):6204--6221, August 2026.
\newblock doi:10.1109/TIT.2026.3686267.
\newblock Also arXiv:2501.12124v4 (last revised 19 August 2025).

\bibitem{etzion-part1}
T.~Etzion,
\newblock On de Bruijn array codes---Part I: Nonlinear codes,
\newblock \emph{IEEE Transactions on Information Theory},
71(2):1434--1449, 2025.
\newblock doi:10.1109/TIT.2024.3518434.

\end{thebibliography}

\end{document}
